% 替换位置：第7.4节起至第7章结束；整体替换原7.4--7.6
% 这是修改片段，不是完整文档；不要整份直接追加到原文末尾。
\subsection{候选生成与联合搜索方法}
\subsubsection{空间候选与服务结构}
采用水平网格、服务区位置和地形高点生成中继候选，离地高度取150米和300米，
经DEM范围、高度及回传链路筛选后形成597个候选位置。
时间候选组合涉及其中49个位置和1904组三架次结构。
求解采用候选筛选与通信服务日历相结合的方法，检查中继服务窗口与运输轨迹的时间匹配。

\subsubsection{通信服务日历参与运输调度}
进一步采用一架中继持续服务、另一架先后执行两个中继任务的结构，
将服务窗口、返航、机体周转、再次准备与建链时间纳入运输候选的可执行性检查。
以运输可行候选作为搜索起点，调整运输机型、批次顺序、访问顺序及准备时刻，
使需要中继的航段避开中继尚未建链或换场的时段。
三处位置沿用候选130、511和489，最终按实际通信需求重新确定服务起止时刻。
运输时刻与中继服务窗口共同参与候选筛选。

\subsubsection{计算限制}
服务日历搜索种子为20260928，时间预算180秒，包含最终细检查的实际耗时约182.748秒。
候选运输批次最多涉及4个服务区；中继位置与三架次结构仍受限制。
本研究服务日历搜索暂要求所有已给定期望时限均满足，以寻找$D=0$方案，
这只是搜索限制，普通物资期望时间在原模型中仍为软目标。
未穷举连续位置、任意组批和中继任务数量，也未证明全部离散候选范围内的最优性。

\subsection{联合调度结果与权衡分析}
\subsubsection{中继位置与任务安排}
所得方案使用2架中继机、3组初始满电能源组件执行3个中继架次。
悬停位置、任务时刻和能耗分别列于下表，所有位置均满足DEM范围及离地高度约束。
\begin{table}[htbp]
\centering
\small
\caption{所得方案的中继悬停位置}
\label{tab:q3_locations}
\begin{tabular}{llrrr}
\toprule
架次 & 机体 & 经度/度 & 纬度/度 & 海拔/m \\
\midrule
R001 & R01 & 109.204167 & 23.048056 & 837.268860 \\
R002 & R02 & 109.270134 & 23.015271 & 627.188568 \\
R003 & R02 & 109.266111 & 23.056389 & 623.466492 \\
\bottomrule
\end{tabular}
\end{table}
\begin{table}[htbp]
\centering
\small
\caption{所得方案的中继任务安排}
\label{tab:q3_relay_schedule}
\begin{tabular}{lrrrrr}
\toprule
架次 & 准备/s & 建链完成/s & 服务结束/s & 返航/s & 能耗/kWh \\
\midrule
R001 & 325.265 & 1056.880 & 5721.485 & 6302.231 & 1.717471 \\
R002 & 130.908 & 738.785 & 3329.431 & 3768.932 & 1.019583 \\
R003 & 4315.907 & 5077.631 & 5506.123 & 6099.161 & 0.457479 \\
\bottomrule
\end{tabular}
\end{table}

R001、R002、R003的返航SOC分别为$46.329029\%$、$68.138039\%$和$85.703779\%$。
R002于$3768.932416\,\mathrm{s}$返航，完成300秒周转后于$4068.932416\,\mathrm{s}$恢复机体可用，
早于R003的准备时刻$4315.906989\,\mathrm{s}$，两次机体占用不冲突。
但R002的能源组件到$4683.137907\,\mathrm{s}$才充满，不能直接承接R003。

\subsubsection{联合方案汇总}
80箱全部按期交付，医疗及首批硬时限全部满足，$D=0$。
\begin{table}[htbp]
\centering
\small
\caption{第三问联合调度结果}
\label{tab:q3_summary}
\begin{tabular}{lr}
\toprule
指标 & 结果 \\
\midrule
运输/中继架次 & 24/3 \\
运输最晚返航/s & 6208.684865 \\
联合最晚返航/s & 6302.230771 \\
运输能耗/kWh & 69.240573590 \\
中继能耗/kWh & 3.194532893 \\
总能耗/kWh & 72.435106483 \\
运输机/运输电池 & 8/14 \\
中继机/能源组件 & 2/3 \\
\bottomrule
\end{tabular}
\end{table}

联合完成时间为$105.037180\,\mathrm{min}$，由R001的最终返航决定，
比运输最晚返航晚$93.545906\,\mathrm{s}$。
第二问与第三问选定的运输子方案具有相同的路线、机型和任务时刻。
该运输子方案满足第二问的运输约束；第三问还需配置相应中继任务，
并独立检查全过程通信，才能形成联合可行方案。

\subsubsection{目标优先级比较}
按$D\rightarrow C_{\mathrm{joint}}\rightarrow E_{\mathrm{total}}\rightarrow N_{\mathrm{total}}$
与$C_{\mathrm{joint}}\rightarrow D\rightarrow E_{\mathrm{total}}\rightarrow N_{\mathrm{total}}$
筛选当前候选池，均选中上述$D=0$、$6302.230771\,\mathrm{s}$、
$72.435106483\,\mathrm{kWh}$及24运输加3中继架次的方案。
这不证明两个完整优化问题具有相同的全局最优解，尤其本研究服务日历搜索偏重全部按期交付的候选。

\subsection{独立轨迹回放与通信验证}
\subsubsection{任务与资源检查}
独立程序重读原始货箱、设备、通信参数和DEM，重构运输全过程轨迹，复算中继飞行与能耗，
检查所有机体、运输电池及中继组件的占用、周转和充电。
两种优先级选中的方案均通过检查，运输最低返航SOC为$20.958111\%$，
最小硬时限余量为$194.866445\,\mathrm{s}$。

\subsubsection{分辨率加密检查}
检查覆盖爬升、巡航、下降、交接以及中继建链与服务边界，结果见表~\ref{tab:q3_refinement}。
\begin{table}[htbp]
\centering
\small
\caption{所得方案的通信分辨率核验}
\label{tab:q3_refinement}
\begin{tabular}{lrrr}
\toprule
时间步长与空间检查 & 点数 & 失败点 & 最小裕量/dB \\
\midrule
1秒，5米射线 & 36218 & 0 & 0.000703269 \\
0.5秒，2.5米射线 & 72113 & 0 & 0.000356918 \\
0.5秒及局部加密，全穿越像元 & 73220 & 0 & 0.000155292 \\
\bottomrule
\end{tabular}
\end{table}

对端点裕量低于$0.02\,\mathrm{dB}$或保障源变化的区间细化至不超过0.05秒，
对定位到的切换点二分至$10^{-7}\,\mathrm{s}$。
最差所选链路裕量出现在P010巡航阶段$t=3860.254957\,\mathrm{s}$附近。

\subsubsection{验证结论与适用范围}
准确结论为：\textbf{在所述时空分辨率下未发现中断}。
离散检查未严格排除采样间隙中的短暂失效，不能写成已数学证明中断时长为零。
全部196个运输阶段保守列入尚未严格认证区间；248条通信保障记录及其任务关联用于第四问冻结。
第一层$D=0$达到非负下界，后续目标无全局最优证明。
标准航程能量标定、爬升效率解释、未计运输交接悬停能耗，
以及中继多机并发与零额外切换耗时仍是需要用户核验的工作假设。
